% Ad Atticum 14.1 --- headnote
% ad-atticum-14-01, 7 April 44 BC, suburban villa of Matius

\textit{Cicero to Atticus, written on 7 April 44 BC at the suburban
villa of Gaius Matius outside Rome --- Perseus dateline
\textit{Scr. in suburbano Mati vii Id. Apr. a. 710 (44)}. This is
the earliest surviving private letter from Cicero after the
assassination, and its first section is the historic page: Cicero
has just called on Matius, the cultivated Caesarian friend of
both men, and reports the verdict back to Atticus in three
hammered phrases --- \textit{nihil perditius; explicari rem non
posse} --- nothing more lost; the thing cannot be untangled. ``If
a mind like his could find no way out, who will find one now?''
Matius is grieved; he predicts trouble in Gaul within twenty
days; he says nothing matters can go on this way. Cicero's
notice of Oppius at the end (the other great Caesarian
intimate, who ``says nothing that could offend any decent
man'') is the contrast that gives Matius's gloom its weight.}

\textit{The second section turns to news from Rome --- Sextus
Pompeius's movements, and above all Marcus Brutus, on whose
political competence Cicero is already half-suspicious. He
copies back to Atticus what Matius reported of Caesar's own
view of Brutus (``whatever he wants, he wants intensely''),
together with an overheard remark of Caesar's about Cicero
himself --- caught waiting in the antechamber, Caesar joked
about Cicero's evident dislike of him. The register is the
intimate, hurried one of the post-Ides retreat: short paratactic
sentences, anecdotes set down \textit{ut enim quidque succurrit
libet scribere} (``for I might as well write whatever comes to
mind''), and a closing that already sounds like the rhythm of
April: \textit{equidem nihil intermittam} --- I shall not let
up.}
